% ===================================================================
%  CADRO N.º 1 — A escola. — O instituto. — A aula.
%
%  Ceci n'est PAS une transcription : c'est une traduction, faite en
%  2026, en galicien normatif (normativa da Real Academia Galega).
%  D'ou trois differences avec texto/io et texto/fr, et trois
%  seulement :
%
%   * pas de \nl ni de \cc ;
%   * pas de \begin{VUpage} ;
%   * le gras se pose sur le TERME seul, ponctuation dehors.
%
%  Les cles « %%K » sont celles de texto/io, macro pour macro pour
%  l'apparat, et les renvois \textsuperscript{(n)} visent les MEMES
%  objets de la planche.
%
%  DOUZIEME COLONNE DES DIX-SEPT, et la quatrieme langue romane apres
%  le catalan, l'occitan et le roumain — mais la premiere qui ait DEUX
%  voisines immediates dans le volume : le portugais et l'espagnol y
%  sont deja, et le galicien se tient entre eux.
%
%  CELA REND LE CONTROLE FACILE ET LA COPIE DANGEREUSE. La regle est
%  celle de toute la serie : la colonne voisine est un temoin, jamais
%  un modele. Ecrire le galicien en partant du portugais donnerait un
%  portugais mal orthographie ; en partant de l'espagnol, un espagnol
%  deguise. On ecrit donc en galicien et l'on VERIFIE sur les deux, ce
%  qui est l'inverse.
%
%  LES TRAITS QUI LE SEPARENT DE SES DEUX VOISINES, et qu'il faut
%  tenir a chaque ligne :
%
%   * le X pour la chuintante — « xiz » la craie, « exercicio »,
%     « xeometría », « xullo », la ou l'espagnol met « ti » ou « j »
%     et le portugais « gi » ou « j » ;
%   * le « -ción » comme l'espagnol, contre le « -ção » portugais ;
%   * PAS DE TEMPS COMPOSES : « rematou » et non « ha rematado » ni
%     « tem acabado » — le galicien n'a pas d'auxiliaire « haber » au
%     parfait, et c'est ce qui l'ecarte le plus de l'espagnol dans la
%     prose ordinaire ;
%   * l'article qui se soude a la preposition : « do », « da »,
%     « no », « na », « ao », « polo », « co » ;
%   * l'infinitif conjugue, qu'il partage avec le portugais et que
%     l'espagnol n'a pas.
%
%  LE CABLAGE, LUI, NE DEMANDE PRESQUE RIEN, et c'est le premier cas
%  depuis l'esperanto. Ses quatre ordinaux feminins sont deja dans la
%  liste — « primeira » et « terceira » par le portugais, « segunda »
%  et « cuarta » par l'espagnol — sa serie s'ecrit « serie » et sa
%  scene « escena », l'une et l'autre deja presentes. Il n'y faut donc
%  ajouter qu'un seul membre, « CADRO » dans NUMERO_TAB, qui est le
%  « CUADRO » espagnol moins son U.
% ===================================================================

%%K t01-apar-0 apar
\VUsaut{2.0mm}
\VUcentre{12.6pt}{120}{LIBRIÑO EXPLICATIVO}
\VUsaut{3.2mm}
\VUcentre{8.4pt}{160}{DOS}
\VUsaut{2.6mm}
\VUtitre{15.6pt}{0}{cadros auxiliares Delmas}
\VUsaut{3.4mm}
\VUornamento{9pt}
\VUsaut{5.0mm}
\VUcentre{13.2pt}{90}{PRIMEIRA SERIE}
\VUsaut{3.0mm}
\VUfilet{16mm}
\VUsaut{5.6mm}

%%K t01-apar-1 apar
\VUtitre{15.0pt}{60}{CADRO N.º 1}
\VUsaut{1.4mm}
\VUtitre{11.4pt}{0}{A escola. --- O instituto. --- A aula.}
\VUsaut{2.2mm}
\VUfilet{20mm}
\VUsaut{6.4mm}

%%K t01-tit-1 sub
\VUpk{10.2pt}{40}{Descrición xeral.}
\VUsaut{3.0mm}

%%K t01-01-1 p
1. --- Este cadro representa unha \VUgras{aula} nun \VUgras{instituto}. Nesta aula
hai oito \VUgras{mesas} \textsuperscript{(18)} e oito \VUgras{bancos}
\textsuperscript{(32)}. En cada banco sentan tres alumnos. O \VUgras{mestre}
\textsuperscript{(7)} senta nunha \VUgras{cadeira} \textsuperscript{(8)} na súa
\VUgras{tarima} \textsuperscript{(6)}.

%%K t01-02-1 p
2. --- Á esquerda da tarima hai un \VUgras{armario} \textsuperscript{(15)}
\VUgras{con vidros.} Á dereita está o \VUgras{encerado} \textsuperscript{(1)} co seu
\VUgras{borrador} \textsuperscript{(2)} e o seu \VUgras{xiz} \textsuperscript{(3)}.

%%K t01-02-2 p
Por riba da tarima penduran \VUgras{cadros murais} \textsuperscript{(9, 11, 12)} e un
\VUgras{mapa} \textsuperscript{(10)}.

%%K t01-03-1 p
3. --- Non todos os alumnos están no seu \VUgras{sitio} \textsuperscript{(17)}. Xoán
\textsuperscript{(4)} está de pé xunto ao encerado; ten \VUgras{o} borrador na man e
borra as \VUgras{figuras xeométricas}.

%%K t01-03-2 p
Pedro está preto da tarima; recolle os \VUgras{exercicios escritos}
\textsuperscript{(5)} e lévallos ao mestre.

%%K t01-04-1 p
4. --- \VUgras{Este} mestre é o mestre de \VUgras{linguas vivas}. Ensínalles aos
alumnos a falar, ler e escribir linguas estranxeiras \VUgras{(alemán, inglés,
español, italiano, ruso} ou \VUgras{francés)}.

%%K t01-05-1 p
5. --- A aula é moi ampla e moi clara. Ao fondo hai unha \VUgras{ventá} larga con
dous \VUgras{postigos} \textsuperscript{(78)} e unha \VUgras{porta de vidro,} para
airear e alumar a aula.

%%K t01-05-2 p
Esta aula non é fría. No medio, para quentala, hai unha \VUgras{estufa} moi boa
\textsuperscript{(46)} co seu \VUgras{tubo} \textsuperscript{(45)}.

%%K t01-05-3 p
Na parede están fixados \VUgras{colgadoiros} \textsuperscript{(58)}; deles penduran
os gorros e os abrigos dos escolares.

%%K t01-tit-2 sub
\VUsaut{5.2mm}
\VUpk{10.2pt}{40}{O patio de xogo.}
\VUsaut{3.6mm}

%%K t01-06-1 p
6. --- O \VUgras{reloxo} \textsuperscript{(63)} marca as nove e cinco minutos.

%%K t01-06-2 p
O \VUgras{recreo} \textsuperscript{(76)} acaba de rematar e a lección comeza de
novo. O recreo faise no \VUgras{patio} \textsuperscript{(75)}. Vemos algúns alumnos
que xogan. Un \VUgras{vixiante} \textsuperscript{(74)} vixíaos.

%%K t01-07-1 p
7. --- No patio vemos ademais persoas diversas, a saber, o \VUgras{inspector}
\textsuperscript{(73)}, o director \VUgras{(o reitor)} \textsuperscript{(71)} e o
\VUgras{censor} \textsuperscript{(72)}.

%%K t01-07-2 p
O inspector inspecciona as escolas, o reitor dirixe o instituto, o censor vixía os
estudos.

%%K t01-07-3 p
A cuarta persoa é o \VUgras{porteiro} \textsuperscript{(77)}. Leva baixo o brazo un
rexistro para anotar as ausencias.

%%K t01-08-1 p
8. --- Ao fondo do patio están os \VUgras{aparellos de ximnasia}
\textsuperscript{(70)} e o taller dos \VUgras{traballos manuais}
\textsuperscript{(69)}.

%%K t01-08-2 p
Á dereita do cadro entrevemos as \VUgras{aulas} de \VUgras{música} e de
\VUgras{debuxo}.

%%K t01-noto-1 noto
\VUnotes{143.2mm}{%
(*) Para os \textit{nomes de bautismo} aconséllase tamén transcribilos segundo a
forma latina. Ese é o único xeito de fuxir de variantes sen conta. Ademais, como
escoller entre \textit{Giovanni, Jean, Johann, Jan, Hans, John, Ivan}, etc.? Imos
crear un dicionario especial para os nomes de bautismo? Tomemos do latín o seu
nominativo e digamos: \VUgras{Ioannes, Ioanna; Iulius, Iulia; Iakobus; Andreas;
Lukas.} (Gram. Completa).%
}

%%K t01-tit-3 sub
\VUpk{10.2pt}{40}{Descrición detallada.}
\VUsaut{2.4mm}

%%K t01-tit-4 sub
\VUpk{10.2pt}{40}{Os bos alumnos.}
\VUsaut{3.0mm}

%%K t01-09-1 p
9. --- Os alumnos do primeiro banco á dereita da tarima son \VUgras{Henrique} (*),
\VUgras{Estevo} e \VUgras{Carlos.}

%%K t01-09-2 p
Henrique é moi atento e traballador; escoita o mestre. Na súa mesa ten un
\VUgras{caderno} \textsuperscript{(28)}, un \VUgras{libro} \textsuperscript{(29)}, un
\VUgras{dicionario} \textsuperscript{(30)} e un \VUgras{estoxo}
\textsuperscript{(31)}. O seu libro ten moitas \VUgras{páxinas}; cada capítulo contén
varios \VUgras{parágrafos}, e cada \VUgras{liña} moitas \VUgras{palabras.}

%%K t01-09-4 p
O seu veciño Estevo \textsuperscript{(24)} é fervoroso. Apunta no seu caderno a
lección para a vez seguinte. Ten unha bela \VUgras{letra.} O seu libro está pechado.
Só vemos a \VUgras{tapa} \textsuperscript{(27)}. Xunto a el está o seu
\VUgras{compás} \textsuperscript{(26)}.

%%K t01-09-5 p
Carlos \textsuperscript{(23)} ten o libro na man; ábreo para reler a súa lección.
Ten, coma cada alumno, un \VUgras{tinteiro} \textsuperscript{(25)} diante de si. É
obediente.

%%K t01-10-1 p
10. --- Na mesa do lado están \VUgras{Lois, Antón} e \VUgras{Paulo.} Lois recita a
súa \VUgras{lección} \textsuperscript{(22)} e responde ás \VUgras{preguntas} do
mestre. Antón \textsuperscript{(21)} é moi desatento; ás veces o mestre mesmo o
castiga. Paulo \textsuperscript{(20)} é bo \VUgras{compañeiro,} amable e servizal. É
moi atento.

%%K t01-tit-5 sub
\VUsaut{4.2mm}
\VUpk{10.2pt}{40}{Os malos alumnos.}
\VUsaut{3.2mm}

%%K t01-11-1 p
11. --- Detrás de Henrique está \VUgras{Xurxo} \textsuperscript{(38)}. Non é mal
rapaz, pero é \VUgras{falador,} desatento e inquedo. A súa \VUgras{lixeireza} e a
súa \VUgras{desobediencia} causan \VUgras{desorde} durante a lección, e a miúdo
merece unha severa \VUgras{advertencia.} O seu \VUgras{caderno de borrancho}
\textsuperscript{(35)} está sucio; \VUgras{mánchao de tinta} coa súa \VUgras{pluma}
\textsuperscript{(34)}, e por iso emprega sempre a súa \VUgras{goma de borrar}
\textsuperscript{(36)}; o seu \VUgras{cartafol} \textsuperscript{(33)} está rachado,
de descoidado que é. Por que trae a súa \VUgras{estilográfica}
\textsuperscript{(37)} á aula \VUgras{das} linguas vivas?

%%K t01-11-2 p
O seu veciño Edmundo \textsuperscript{(39)} non o quere. É certo que é algo
indolente, pero é calado e tranquilo. Non anda a falar; mais, en vez de escoitar, prefire
ler os \VUgras{contos} do seu \VUgras{libro de lectura}.

%%K t01-11-3 p
O terceiro alumno deste banco é \VUgras{Ludovico} \textsuperscript{(40)}. É
diligente. Escribe nunha \VUgras{folla de papel} branca. Diante de si puxo o
\VUgras{papel secante} \textsuperscript{(41)}.

%%K t01-12-1 p
12. --- Os alumnos do segundo banco, á esquerda do mestre, son moi novos. O primeiro,
o pequeno \VUgras{Emilio,} aínda non sabe escribir moi ben; trae a súa
\VUgras{lousa} \textsuperscript{(42)}, o seu \VUgras{lapis de lousa} e a súa
\VUgras{esponxa} \textsuperscript{(43)}.

%%K t01-12-2 p
Á beira del están \VUgras{Augusto} e \VUgras{Bernardo} \textsuperscript{(44)}. Este
último traballa ben, pero a súa \VUgras{roupa} está moi descoidada.

%%K t01-13-1 p
13. --- Detrás deles están tres \VUgras{preguiceiros. Alberte} non aprende as súas
leccións. \VUgras{Xacobe} \textsuperscript{(47)} dorme en vez de escoitar. A súa
\VUgras{preguiza} vállelle \VUgras{reproches} e mesmo \VUgras{castigos.} Finalmente
\VUgras{Cristián} \textsuperscript{(48)} é lixeiro de máis.
\VUgras{Pórtase} mal e non fai ningún esforzo por emendarse.

%%K t01-14-1 p
14. --- Os seus veciños, pola contra, pórtanse ben. \VUgras{Ernesto}
\textsuperscript{(51)} gusta da orde e da regularidade; emprega sempre a súa
\VUgras{regra} \textsuperscript{(50)} e o seu \VUgras{lapis} \textsuperscript{(49)}.
É \VUgras{alumno externo.}

%%K t01-14-2 p
\VUgras{Francisco} \textsuperscript{(52)} é \VUgras{interno}; viste uniforme, come e
dorme no instituto. É bo \VUgras{compañeiro.}

%%K t01-14-3 p
Mauricio afía o seu \VUgras{lapis} coa súa \VUgras{navalla}
\textsuperscript{(56)}; é moi coidadoso.

%%K t01-14-4 p
Trae todos os seus libros, a súa \VUgras{gramática} \textsuperscript{(55)}, o seu
\VUgras{caderniño de notas} \textsuperscript{(54)} e mesmo un \VUgras{cartapacio}
\textsuperscript{(53)}. A súa \VUgras{mochila} \textsuperscript{(57)} está no chan
detrás del.

%%K t01-14-5 p
As dúas mesas do fondo da aula están ocupadas polos \VUgras{bos alumnos}
\textsuperscript{(19)}.

%%K t01-tit-6 sub
\VUsaut{4.6mm}
\VUpk{10.2pt}{40}{Debuxo e música.}
\VUsaut{3.4mm}

%%K t01-15-1 p
15. --- Na \VUgras{aula de debuxo} hai \VUgras{modelos} graciosos pendurados na
parede e unha \VUgras{lámpada} \textsuperscript{(62)} coa súa \VUgras{pantalla,}
colgada do teito. O \VUgras{mestre} \textsuperscript{(60)} é moi paciente; corrixe os
\VUgras{debuxos} \textsuperscript{(61)} dos alumnos e ensínalles a debuxar do natural
un \VUgras{busto} \textsuperscript{(59)}.

%%K t01-16-1 p
16. --- A \VUgras{aula de música} semella moi interesante. Nela apréndese a ler a
\VUgras{música}, a solfexar as \VUgras{notas da escala} \textsuperscript{(67)}
escritas no \VUgras{pentagrama} (do, re, mi, fa, sol, la,
si\,=\,C. D. E. F. G. A. B.), a coñecer as \VUgras{claves} e a cantar
\VUgras{melodías} fermosas. Póñense as partituras no \VUgras{atril}
\textsuperscript{(65)}. Un dos alumnos \VUgras{toca o piano}
\textsuperscript{(64)} e o \VUgras{mestre de música} \textsuperscript{(66)}
\VUgras{marca o compás} \textsuperscript{(68)}.

%%K t01-17-1 p
17. --- Entrevemos un recanto do taller dos \VUgras{traballos manuais}
\textsuperscript{(69)}, onde os rapaces poden aprender a cepillar a madeira e a limar
o ferro. Iso non existe en todos os institutos.

%%K t01-tit-7 sub
\VUsaut{5.0mm}
\VUpk{10.2pt}{40}{A aula das ciencias.}
\VUsaut{3.6mm}

%%K t01-18-1 p
18. --- O mestre de matemáticas ensínalles \VUgras{aos} alumnos
\VUgras{xeometría, aritmética, cálculo} e o \VUgras{sistema métrico}. Vemos no cadro
as principais figuras xeométricas: o \VUgras{círculo} \textsuperscript{(a)}, o
\VUgras{cadrado} \textsuperscript{(b)}, o \VUgras{triángulo} \textsuperscript{(c)}, a
\VUgras{liña} \textsuperscript{(d)}.

%%K t01-18-2 p
Vemos tamén unha \VUgras{operación}; non é unha \VUgras{suma,} nin unha
\VUgras{resta,} nin unha \VUgras{multiplicación}: é unha \VUgras{división}
\textsuperscript{(e)}.

%%K t01-19-1 p
19. --- No cadro do sistema métrico vemos: o \VUgras{metro} \textsuperscript{(a)}, a
\VUgras{balanza} \textsuperscript{(b)} e as súas \VUgras{pesas,} o \VUgras{litro}
\textsuperscript{(e)}, o \VUgras{decalitro} \textsuperscript{(f)}, o
\VUgras{decilitro} e o \VUgras{metro cúbico} \textsuperscript{(d)}.

%%K t01-19-2 p
Os alumnos estudan tamén \VUgras{ciencias naturais} \textsuperscript{(11)},
zooloxía \textsuperscript{(a)}, botánica \textsuperscript{(b)}, xeoloxía
\textsuperscript{(c)}, e a \VUgras{historia} \textsuperscript{(12)}.

%%K t01-19-3 p
No armario están os aparellos de \VUgras{física} \textsuperscript{(15)} e de
\VUgras{química} \textsuperscript{(16)}.

%%K t01-19-4 p
Enriba do armario están: a \VUgras{esfera} \textsuperscript{(14)}, o \VUgras{cono} e
o \VUgras{globo terráqueo} \textsuperscript{(13)}.

%%K t01-19-5 p
Por riba dos cadros pendura un \VUgras{mapa} que representa as cinco partes do mundo:
\VUgras{Europa} \textsuperscript{(g)}, \VUgras{Asia} \textsuperscript{(e)},
\VUgras{África} \textsuperscript{(f)}, \VUgras{América} \textsuperscript{(ab)} e
\VUgras{Oceanía} \textsuperscript{(h)}, e os dous grandes océanos, o
\VUgras{Pacífico} \textsuperscript{(c)} e o \VUgras{Atlántico} \textsuperscript{(d)}.
\VUsaut{6.0mm}
\VUfilet{22mm}
